\documentclass[10pt,a4paper]{amsart}
\usepackage[margin=19mm]{geometry}
\usepackage{amsmath,amssymb,mathtools}
\usepackage[scr=rsfs,cal=euler]{mathalfa}
\usepackage{xfrac}
\usepackage[unicode,hidelinks,bookmarks=false]{hyperref}
\newcommand{\ie}{i.e.\ }
\newcommand{\cf}{cf.\ }
\newcommand{\resp}{respectively\ }
\newcommand{\Subsection}{\subsection}
\providecommand{\nobreakdash}{\nobreak\hbox{-}}
\setlength{\emergencystretch}{2em}
\multlinegap0pt
\usepackage[shortlabels]{enumitem}
\usepackage{amssymb}
\usepackage[normalem]{ulem}
\usepackage{tikz}
\newtheorem{prop}{Proposition}
\newtheorem{observation}{Observation}
\newcommand \Lin {\mathfrak{L}}
\newcommand \Su {\mathscr{S}}
\newcommand \Ult {\mathfrak{U}}
\newcommand \dom {\operatorname{dom}}
\newcommand \fin {{\operatorname{fin}}}
\newcommand \rng {\operatorname{rng}}
\newcommand \seq [1]{\langle #1 \rangle}
\newcommand \set [1]{\{#1\}}
\title{Generic dc-automorphisms of two-sorted ultrametric spaces}
\author{Adam Barto\v{s}, Wies{\l}aw Kubi\'s, Aleksandra Kwiatkowska, Maciej Malicki}
\date{Source: arXiv:2606.11498v1 (CC BY 4.0)}
\begin{document}
\maketitle

\noindent\emph{Source and licence.} Generic dc-automorphisms of two-sorted ultrametric spaces. Version arXiv:2606.11498v1 (CC BY 4.0), distributed by arXiv under the Creative Commons Attribution 4.0 International licence, http://creativecommons.org/licenses/by/4.0/, as recorded in the arXiv licence metadata for this exact version and shown on the abstract page https://arxiv.org/abs/2606.11498v1. Full licence notice: \texttt{licenses/arxiv-2606.11498-CC-BY-4.0.txt}.\par\smallskip

\noindent\textit{Editorial scope.} Counted unit: the article's prop thm:unique\_onepoint\_extension (source0.tex lines 1545--1549). The article's complete original proof of it is delivered with it (source0.tex lines 1551--1576, one \texttt{proof} environment of 26 lines). No mathematics is written, repaired or re-derived: the statement and the proof are the article's own lines, in the article's own notation, and no retained line is substituted or re-flowed.  Retained uncounted as the source-local context the proof needs: lines 1285--1292; lines 1537--1543.  Nothing was omitted from any retained span, so the package carries no omission record.  No retained line contains a citation, so the wrapper carries no bibliography and no bibliography entry.  No retained line contains a figure environment, an image-inclusion command or drawing code, so no image is delivered and the package declares no asset dependency.  Editorial formatting only, in addition: an AMS \texttt{amsart} preamble that restates the article's own declarations and macros used by the retained lines, namely the environments \texttt{prop}, \texttt{observation} on separate counters, together with the standard packages those lines need.  The retained environments print in this document as Proposition 1, Observation 1 in order of appearance, whereas the article prints the same items with the numbering of its own full document.  The complete native source of the article as distributed in the arXiv e-print for this exact version is bundled unchanged as \texttt{sources/202555/source0.tex}.
\begin{prop}\label{lo_totaliz}
    $(\Lin^0)^*_\fin$  has CAP and for $\seq{A,p}\in (\Lin^0)^*_\fin$ the following conditions are equivalent:
    \begin{enumerate}[label={\rm(\alph*)}]
        \item $\seq{A,p}$ is an amalgamation base;
        \item $\seq{A,p}$ is determined;  
        \item every orbit of $\seq{A,p}$ is stable.
    \end{enumerate}
\end{prop}

\begin{observation} \label{distance_sort_stable}
    We have that $\seq{A, p} \in \Ult^*_\fin$ is in $\Su$ if and only if every orbit of the point sort is stable and $\seq{D_A, D_p}$ is an amalgamation base in $(\Lin^0)^*_\fin$ (or equivalently is determined in $(\Lin^0)^*_\fin$).
\end{observation}
\begin{proof}
    The claim follows immediately from Proposition~\ref{lo_totaliz} once we observe that an orbit of $\seq{A, p}$ of the distance sort is stable in $\Ult^*_\fin$ if and only if it is stable in $(\Lin^0)^*_\fin$ as an orbit of $\seq{D_A, D_p}$, i.e. we can close or merge an orbit of $D_p$ in an extension $\seq{E, \phi} \geq \seq{D_A, D_p}$ if and only if we can close or merge the orbit in an extension $\seq{B, q} \geq \seq{A, p}$.
    The last equivalence is true because every such extension $\seq{B, q} \geq \seq{A, p}$ induces $\seq{E, \phi} = \seq{D_B, D_q} \geq \seq{D_A, D_p}$, and every extension $\seq{E, \phi} \geq \seq{D_A, D_p}$ induces $\seq{A, p, E, \phi} \geq \seq{A, p, D_A, D_p}$, and we let $\seq{B, q, D_B, D_q} = \seq{A, p, E, \phi}$.
\end{proof}

\begin{prop} \label{thm:unique_onepoint_extension}
    Let $\seq{A, p} \in \Su$.
    For every $x \in A \setminus \dom(p)$ there is a unique (up to isomorphism) forward one-point orbit extension $\seq{A', p'}$ at $x$.
    Moreover, every $r \in D_{A'} \setminus D_A$ is in an extension of an orbit of $D_p$.
\end{prop}

\begin{proof}
    Let $\seq{A', p'} = \seq{A \cup \set{y}, p \cup \set{x \mapsto y}}$ be any such extension.
    It is enough to show that for every $a \in A$ the distance $d_{A'}(a, y)$ is determined by the distance part $D_{p'}$ and by an old distance.
    Specifically, $d_{A'}(a, y) = \max\set{D_{p'}(r_a), s_a}$, where $r_a, s_a \in D_A$ depend only on $a$ and $x$.
    This is indeed enough since the extension $D_{p'} \geq D_p$ is unique up to isomorphism as $D_p$ is determined by Observation~\ref{distance_sort_stable}.

    Let $a \in A$.
    If $a \in \rng(p)$ and so $a = p(a')$ for some $a' \in A$, we have $d(a, y) = D_{p'}(d(a', x))$, and hence in the formula above we can  put $r_a = d(a', x)$ and $s_a = 0$.
    
    Otherwise, $a \notin \rng(p)$.
    Let $b \in \dom(p)$ be such that the distance $d(p(b), a)$ is minimal possible.
    Then we define $r_a = d(b, x)$ and $s_a = d(p(b), a)$.
    We consider the triangle $y, a, p(b)$ with distances $d(y, a)$, $s_a$, and $d(y, p(b))= D_{p'}(r_a)$.
    Clearly, if $s_a \neq D_{p'}(r_a)$, then $d(y, a) = \max\set{s_a, D_{p'}(r_a)}$ and we are done.
    Otherwise, we have $d(y, a) = s_a = D_{p'}(r_a)$ and we are also done, or $d(y, a) < s_a = D_{p'}(r_a)$.
    We claim the last case contradicts our assumptions as then $q = p \cup \set{x \mapsto a}$ induces a partial automorphism of $A$.
    For that it is enough to show that $\phi(d(x, t)) = d(a, p(t))$ for every $t \in \dom(p)$, where $\phi$ is the unique totalization of $D_{p'}$ in $(\Lin^0)^*$, as then letting $D_q$ be the restriction of $\phi$ to $d[\dom(q)^2]$ turns $q$ into a partial dc-automorphism.
    Let $t \in \dom(p)$.
    We have $d_{A'}(y, a) < s_a = d(p(b), a) \leq d(p(t), a)$.
    And therefore by considering the triangle $y, a, p(t)$ in $A'$ we have
    \[
        d(a, p(t)) = d_{A'}(y, p(t)) = D_{p'}(d(x, t)) = \phi(d(x, t)) = D_q(d(x, t)).
    \]
%
    The moreover part of the proposition follows from $d_{A'}(a, y) = \max\set{D_{p'}(r_a), s_a}$.
\end{proof}

\end{document}
